% SPDX-License-Identifier: Apache-2.0
% covers: Hart
\datasheet{Hart}{The Vreteno core and its Sv32 unit, joined, with the core's ports and nothing more.}

\dsfacts{\code{cpu/vreteno/src/hart.rs}}{\code{vreteno32::hart::Hart}}%
{\code{//docs:vreteno}}

\subsection*{Function}

\code{Hart<IW>} is a unit of two units: a \code{Vreteno<IW>} core and
the \code{Mmu8} that translates its addresses for Sv32 (issue 1014).
Its ports are the core's, all of them and no others, so everything
outside sees one core: the bus side, the interrupts, the debug
module's lines and the retire outputs.

Inside, the core tells the unit \code{satp}, the mode a data access is
checked as, \code{sum}, \code{mxr} and the flush, and asks on its two
request wires; the unit answers on two answer wires. The walker's page
table reads go to the core on a channel, which the core puts on its
own bus port as loads, and the words read come back to
the unit on a second channel. The hart therefore needs no bus port of
its own for the page tables.

\code{Hart::with} loads a program into the core's instruction memory.

\subsection*{Parameters}

\begin{center}\footnotesize
\begin{tabular}{@{}lp{0.7\textwidth}@{}}
\toprule
Parameter & Meaning \\
\midrule
\code{IW} & The width of the AXI identifier, passed to the core. The
tables use 2, as the board and the runs do. \\
\bottomrule
\end{tabular}
\end{center}

\dstables{Hart}

\subsection*{Behaviour}

\begin{itemize}
\item The unit goes first in the join. Its answers are registers, and
the core reads them in the same step.
\item The core's requests are registers too, and the unit reads them a
step late in simulation and on time in hardware. So the core waits two
cycles on every request before it believes an answer, which is right
in both, since the unit holds its answer while the request is held.
\item The hart adds no register and no logic of its own: every
register in it is the core's or the unit's.
\end{itemize}

\subsection*{Verification}

\begin{itemize}
\item \code{run.rs}'s machine, which \code{//cpu/vreteno:lockstep\_test},
\code{//cpu/vreteno:hello\_test}, \code{//cpu/vreteno:icache\_test} and
the core's other tests use, runs a \code{Hart}. Lockstep holds it
against the reference model after every cycle, with paging directed
and over random tables.
\item The build simulates \code{Hart<2>}'s netlist, under the entity
name \code{vreteno} and with the demonstration program in the core's
instruction memory, against the trace of \code{//cpu/vreteno:vreteno}
under nvc and Verilator: \code{//docs:vreteno\_sim\_vreteno}\allowbreak\code{\_vreteno\_tb\_test}
and \code{//docs:vreteno\_vsim\_vreteno\_test}.
\item \code{//cpu/vreteno:openocd\_test} runs a stock OpenOCD against a
hart on the simulated board, and \code{//cpu/vreteno:pair\_test} runs
two harts in step.
\item \code{//cpu/vreteno:vreteno\_synth} and
\code{//cpu/vreteno:vreteno\_pnr}, which are manual, synthesise and
place the hart's netlist; on the board it is in every flagship.
\end{itemize}

\subsection*{Used by}

\code{Hart<2>} as \code{cpu} in \code{Board}, and so in the flagship;
two of them in \code{Pair}; and the machine of \code{run.rs}.

\subsection*{Limits}

One hart: there is no second core sharing the unit or the bus port.
The core's fetch and data requests each take the unit's cycle of
latency, which the core's own rules account for (\code{Vreteno}'s
sheet).
